%% Finite-tree automata/MSO proof for the displayed ranked-tree signature.

\section{Finite trees and monadic second-order logic}

\begin{annexitem}{annex:p9.trees.1}{thm:p09-finite-tree-automata}{Finite-tree automata and monadic second-order logic}
\begin{proof}
Fix the finite ranked alphabet \(\Sigma\).
The subset invariant in Section~\ref{sec:p09-finite-tree-determinization}
proves determinization, including the empty subset.
The disjoint-union, product, and complete-complement constructions in
Section~\ref{sec:p09-finite-tree-boolean} prove the three Boolean closures.
The upward marking invariant and its occurrence-counter implementation in
Section~\ref{sec:p09-finite-tree-emptiness} prove membership and emptiness
decidability.  We prove the two logical translations explicitly.

\paragraph{Marking assignments on a fixed input tree.}
List \(m\) free individual variables and \(n\) free set variables.
Use the decorated ranked alphabet
\[
 \Sigma^{m,n}
   =\{(f,\mathbf b):f\in\Sigma,\ \mathbf b\in\{0,1\}^{m+n}\},
 \qquad
 \operatorname{arity}(f,\mathbf b)=\operatorname{arity}(f).
\]
The underlying \(\Sigma\)-tree is obtained by forgetting the bits.
An individual-variable track marks exactly its assigned node; a set
track marks exactly the nodes in the assigned subset.
Different individual variables may mark the same node.

The language \(V_{m,n}\) of legal assignment markings is recognizable.
For each individual track, a state records its number of marked nodes
in the current subtree, truncated at two.
At a node, add the current bit and the child counts and truncate
again.  Accept exactly when every individual count is one.
The set-variable tracks impose no extra restriction.
There are finitely many count vectors, and induction on the tree
proves the count invariant.
When \(m=0\), the empty count vector accepts every decorated tree.
When \(\Sigma\) has no ground terms, all these languages are empty,
which is the correct empty class of input structures.

\paragraph{Rank-preserving erasure and extra tracks.}
Let \(h\) erase selected bit coordinates, preserving the ranked symbol
and its arity.  The image under \(h\) of a recognizable decorated-tree
language is recognizable: replace the label of every transition by
its projected label, retaining its child states and target.
An original run projects to a run of this automaton.
Conversely, for each node of a finite projected run choose an
original label witnessing its transition.  Those choices give one
decorated input tree and a run of the original automaton.
There is no compatibility condition between bits at different nodes
beyond what its states and transitions already enforce.

To read an automaton over a larger list of tracks, let each new
transition ignore the extra bits.  This recognizes \(h^{-1}(L)\);
intersecting with the legal-marking language for the larger list
enforces the individual-variable conditions.
Thus automata for two subformulas can always be placed on the same
ordered union of their free-variable tracks.

\paragraph{Atomic formulas.}
Work relative to \(V_{m,n}\).
Each of the following tests is recognized by a finite bottom-up
automaton, then intersected with \(V_{m,n}\).
For equality \(x=y\), record whether some node carries both individual
markers.  Under the singleton condition this is exactly equality
of their assigned nodes.
For a label atom \(P_f(x)\), record whether an \(x\)-marked node
has label \(f\).
For \(x\in X\), record whether a node carries both the \(x\)-marker
and the \(X\)-bit.
In each case the record is a Boolean value: the current local event
or the event already recorded in one of the child subtrees.
Bottom and top are represented by the empty language and \(V_{m,n}\).

For the child relation \(E_i(x,y)\), states additionally remember
the individual-marker bits at the root of the current subtree.
At a node labelled \(f\) of arity \(r\), the new local event holds
exactly when \(i\leq r\), the current node has its \(x\)-bit set,
and the root bits in its \(i\)-th child state contain the \(y\)-marker.
Combine this event with the child event flags.
The newly stored root bits are those of the current label.
Induction shows that the final flag says exactly that the two
uniquely marked nodes satisfy \(E_i\).
If a formula uses the same individual variable twice, it uses the
same track; the construction therefore also covers \(x=x\) and
\(E_i(x,x)\).

\paragraph{Induction on formulas.}
Associate with each monadic second-order formula \(\varphi\) an
automaton accepting exactly its legal marked trees on which
\(\varphi\) is true.
The atomic constructions begin the induction.
For conjunction and disjunction, extend both automata to the same
free-variable track list and take intersection or union.
For negation, use relative complement
\[
 L_{\neg\varphi}=V_{m,n}\setminus L_\varphi.
\]
The legal-marking restriction is part of the invariant, including
after negation.

For \(\exists X\,\varphi\), construct the body automaton with the
additional set track, then erase that track.
Its finitely many node bits range over exactly all subsets of the
finite input domain, so the image operation implements existential
set quantification.
For \(\exists x\,\varphi\), use an additional individual track,
enforce its singleton condition, and then erase it.
Every accepted expansion supplies a node witness, and every node
witness supplies that marking.
If the quantified variable does not occur in the body, add its
unused track first; the nonempty input tree still has a singleton
marking for an individual variable.
Bound variables are renamed to avoid clashes before these operations.
Universal quantifiers are obtained by negation and existential
quantification.  Set equality and inclusion, if used as abbreviations,
are expressed by their pointwise membership formulas.
These cases exhaust the usual monadic second-order syntax in the
stated label/child/equality signature.

For a sentence, no tracks remain.  The alphabet
\(\Sigma^{0,0}\) is canonically \(\Sigma\), and its automaton therefore
recognizes precisely the sentence's language of finite ranked trees.

\paragraph{Describing an accepting run by sets.}
Conversely, let \(\mathcal A=(Q,\Sigma,\Delta,F)\) be a finite
bottom-up tree automaton.
Introduce one monadic set variable \(Z_q\) for each state \(q\in Q\).
Require these sets to partition the input domain:
\[
 \forall x\left[
  \left(\bigvee_{q\in Q}x\in Z_q\right)
  \land
  \bigwedge_{\substack{q,r\in Q\\q\ne r}}
                 \neg(x\in Z_q\land x\in Z_r)
 \right].
\]
Let \(d\) be the largest symbol arity, taking \(d=0\) when
\(\Sigma\) is empty, and define
\[
 \operatorname{Root}(x)
   \equiv\neg\exists y\,\bigvee_{i=1}^{d}E_i(y,x).
\]
On every input tree this picks its unique root.
Require
\[
 \exists x\,
       \left(\operatorname{Root}(x)\land
                         \bigvee_{q\in F}x\in Z_q\right).
\]
For each \(f\in\Sigma\) of arity \(r\), include
\[
 \forall x\left[
 P_f(x)\to
 \bigvee_{\substack{f(q_1,\ldots,q_r)\to q\\\in\Delta}}
 \left(
  x\in Z_q\land
  \exists y_1\cdots y_r
       \bigwedge_{i=1}^{r}
                    (E_i(x,y_i)\land y_i\in Z_{q_i})
 \right)\right].
\]
For \(r=0\), the existential block is empty and its conjunction is
\(\top\), so this is exactly the nullary-transition condition.
An empty disjunction is \(\bot\).
Existentially quantify all the \(Z_q\) around the conjunction of
partition, root, and transition clauses.

An accepting run supplies its state classes as witnesses to these sets.
Conversely, a satisfying partition assigns exactly one state to each
node.  Its actual label selects a transition clause, and its ranked
tree children are precisely the witnesses in that clause.
Thus every node satisfies a legal transition, and the root clause
makes the resulting run accepting.
If \(Q\) or \(F\) is empty, the same empty-disjunction conventions
make the formula unsatisfiable on a nonempty tree, as required.
This proves the converse translation.

All subsets of the finite input domain are finite.
Accordingly weak and full monadic quantification have the same
meaning on these structures.  The two translations and the finite
automaton constructions establish all assertions of the theorem.
\end{proof}
\end{annexitem}
